% !TEX root = ../journal_0421052099.tex
% ---------------------------------------------------------------------
% The instrument, at the width the attribution study needs and no more.
%
% CUT 2026-09-23 (third pass) under a hard 20-22 page budget for
% abstract-to-conclusion. The mechanics MOVED, they were not deleted:
% the tool API, eq:score, the evaluator/router/backtracker, the
% three-route claim path, the budget caps and sec:removal are now
% Appendix B (sections/appendix-instrument.tex), which the budget does
% not count and which this section points at four times. What stays
% here is what Sections 5-8 actually reason from: the loop is a
% program, so a component can be removed by a flag; every path reaches
% the Answerer through the Verifier; the verifier's graph routes are
% relation-blind; and the layer is nearly free.
%
% The relation-blind sentence is load bearing for the introduction and
% for sec:limits-contribution and must not leave the body.
%
% References to sec:budgets, sec:removal and the app: labels are \ref,
% not \Cref: they are appendix subsections now, and elsarticle's
% \appendix redefines \thesection, so \Cref prints "Section Appendix
% B.5" where "Appendix B.5" is wanted.
%
% Earlier passes removed alg:agr, fig:claim-path, tab:budget, the Yale
% worked example and the verify_triple history from the body; the last
% two are in Appendix B, the first three are gone for good (the thesis
% carries them).
%
% NOTE: journal/figures/fig_claim_path.tex is still pinned to the
% thesis's copy by tests/test_paper_self_contained.py though the
% manuscript no longer \inputs it.
% ---------------------------------------------------------------------

\section{The Instrument: AGR}\label{sec:framework}

AGR inverts the usual agent pattern. It does not hand the model a set of
tools through native function-calling and let the model drive the loop.
A program orchestrates every step, calls the graph tools
deterministically, and consults the model only at bounded decision
points, through prompts constrained to return JSON. Every graph
operation is therefore logged by construction rather than reconstructed
from the model's account of what it did, and budgets are enforced by
counters a router checks rather than requested in a prompt the model may
ignore. The consequence this paper uses is that each of the loop's four
components can be removed by a configuration flag without touching the
others, which is what \Cref{sec:attribution} does; the supplementary
material states exactly what each flag removes.

\subsection{State Machine}\label{sec:statemachine}

A state-machine diagram in the supplementary material gives the
topology: six agent nodes over a shared typed state, with three cycles.
The \textbf{Planner} decomposes the question into an ordered chain of at
most four sub-objectives in one model call, from a prompt whose first
exemplar is a single-hop question that is \emph{not} decomposed.
\textbf{Explorer} and \textbf{Evaluator} form a search loop: the
explorer expands the frontier by scoring every candidate relation of at
most five anchors, blending embedding similarity against the active
sub-objective with a batched model judgement at a frozen $\alpha = 0.7$,
and the evaluator judges whether that objective is satisfied from the
thirty most similar traversed triples. It may accept only entities the
graph returned. The \textbf{Backtracker} restores an earlier,
higher-scoring frontier and bans the edges that led away from it; that
ban list is what turns backtracking into a search operation, since
without it a deterministic scorer facing the same restored frontier
would select the same edges again. The supplementary material gives the
four graph operations, the scoring rule, and the router and the
backtracker's snapshot discipline. Every path from Evaluator to Answerer
passes through the \textbf{Verifier}, and the budgets recorded in the
supplementary material bound all three cycles, so no sequence of model
outputs can produce a non-terminating run.

\subsection{Structural Verification Layer}\label{sec:verification}

Nothing in the navigation loop constrains what the final answer
\emph{asserts}. The verifier is the check that does, and it is the
component the system was designed around. One model call drafts prose
from the traversed triples, extracts the entities the draft asserts as
its answer, and decomposes the draft into atomic claims $\langle h, r, t
\rangle$ over entity \emph{names}, since the model has never seen a
graph identifier. Each claim is then routed through up to three tests
and leaves at the first that settles it: \textbf{traversed adjacency},
which asks whether the claim's endpoints were walked together and is the
only route of the three that attaches citable evidence;
\textbf{\textsf{verify\_connection}}, which asks the full graph the same
question for endpoints that were reached but never walked together; and,
for whatever clears neither, an \textbf{entailment check} against the
fact window the drafter saw. Symbols decide, and the model is consulted
only where the graph is silent. If no claim is unsupported the draft is
emitted verbatim; otherwise the verifier either sends the agent back to
explore for the missing support or, when the budget is spent, rewrites
the draft around a deterministically filtered entity list. The
supplementary material gives the routing and the three outcomes in full.

\textbf{What the graph routes actually check bounds every claim made for
this layer.} Both of them ignore the relation and the direction, so what
they establish is that the endpoints of a claim are connected in the
environment, not that the claimed relation is the one the edge records:
a claim that X is Y's mother is certified by an edge recording that X is
Y's child. Only one of the three routes attaches evidence at all.
\Cref{sec:limits-contribution} states the exposure that follows and how
much of it the committed record can bound, and \Cref{sec:nulls} measures
what the layer changes.

The layer is also nearly free, which matters because a cost argument is
sometimes offered in place of an accuracy one. The grounded path adds
nothing beyond the single drafting call already required, and at test
scale removing the layer altogether saves $2\%$ and $6\%$ of tokens on
the two ablation half-splits (\Cref{tab:ablation}). Whatever the case
for or against it, it is not a cost argument.
